\documentclass[11pt,a4paper]{article}
\usepackage{kotex}
\usepackage[margin=1in]{geometry}
\usepackage{booktabs}
\usepackage{hyperref}
\usepackage{array}

\title{Nexus Entertainment 프로젝트 실증분석 보고서}
\author{ABC Project}
\date{2026년 5월 11일}

\begin{document}
\maketitle

\section{분석 목표와 범위}
본 보고서는 Nexus Entertainment 디지털 휴먼 트윈 프로젝트의 현재 구현을 대상으로,
행동 데이터 파이프라인이 실행 가능하고 재현 가능한지, 공개 샘플 데이터와 무작위 반례에서
안정적으로 동작하는지, 그리고 사전에 정의한 플레이 스타일 카테고리를 일관되게 산출하는지를
검증한다. 여기서 분석 대상은 사람의 실제 성격이 아니라 관찰 가능한 플레이 스타일과 행동
카테고리다. 분석은 새로 추가한 \texttt{backend/empirical\_analysis.py},
\texttt{backend/validate\_playstyle\_categories.py}, 기존 테스트/검증 스크립트를 함께 사용했다.

\section{사전 기대치와 이론적 최선}
실험 전 이론적 최선은 다음과 같이 정의했다. 단위 테스트와 프론트엔드 테스트는 100\% 통과,
파서 정확도와 엣지 케이스 처리율은 100\%, 무작위 강건성 탐색은 모든 trial에서 예외 없이
메트릭 불변식을 만족하는 상태를 최선으로 보았다. 이 기준은 실제 결과를 과대평가하지 않기 위한
상한선이며, 데이터셋 규모나 실제 사용자 다양성까지 보장하는 기준은 아니다.

\section{방법론}
\subsection{1단계: 실행 가능성 검증}
먼저 백엔드 pytest와 프론트엔드 Vitest를 실행했다. 초기 실행에서는 백엔드 3개 테스트와
프론트엔드 7개 테스트가 실패했다. 이후 후속 작업에서 planning time 정의 불일치도 추가로
수정했다. 주요 원인은 (1) \texttt{/api/behavior} 응답에서
\texttt{synthetic\_theta}, \texttt{synthetic\_beta}가 누락된 API 계약 회귀,
(2) 스트레스 감지 함수가 평탄화된 세션 메트릭을 읽지 못하는 문제,
(3) \texttt{BehaviorTracker}의 named export/호환 메서드 부재,
(4) fake timer 기반 welcome animation 테스트의 비결정성,
(5) \texttt{planning\_time}이 평가 정의와 다른 준비 행동 내부 span을 반환한 문제다.
해당 결함을 수정한 후 동일 명령을 재실행했다.

\subsection{2단계: 하위 컴포넌트와 반례 탐색}
기존 \texttt{backend/comprehensive\_test.py}를 재실행해 이벤트 파싱 정확도, 성능,
일관성, 엣지 케이스, 공개 데이터 파일 존재 여부를 확인했다. 또한
\texttt{backend/empirical\_analysis.py}에서 고정 seed \texttt{20260511}을 사용해
200개의 무작위 trial을 수행했다. 각 trial에서는 단 하나의 변형만 적용했다:
\texttt{drop\_timestamp}, \texttt{drop\_position}, \texttt{unknown\_type},
\texttt{negative\_timestamp}, \texttt{shuffle\_order}. 성공 기준은 파서가 예외 없이
동작하고, \texttt{complexity}, \texttt{path\_efficiency}, \texttt{risk\_taking},
\texttt{diversity}가 모두 [0, 1] 범위에 있으며, 결과가 \texttt{GameBehaviorProcessor}까지
전달되는 것이다.

\subsection{3단계: 전체 평가와 성능 분포}
\texttt{backend/final\_verification\_test.py}로 API, 게임 파이프라인, 성능, E2E 통합을
검증했다. 성능 분석은 이벤트 수만 바꾸는 단일 변수 실험으로 구성했으며, 각 이벤트 수에 대해
30회 반복 후 median과 p95 latency를 기록했다.

\subsection{플레이 스타일 카테고리 검증}
\texttt{datasets/validation/labeled\_playstyle\_sessions.json}에 synthetic calibration fixture를
추가했다. 이 fixture는 실제 사용자 성격 라벨이 아니라, 검증 스크립트가 사전 정의된 카테고리를
재현 가능하게 산출하는지 확인하기 위한 규칙 기반 예시다. 카테고리는 \texttt{planner},
\texttt{iterative\_refiner}, \texttt{route\_optimizer}, \texttt{complex\_builder},
\texttt{resource\_diversifier}, \texttt{risk\_taker}로 정의했다. 후속 실제 데이터 검증을 위해
\texttt{datasets/validation/real\_playstyle\_sessions\_template.json}와
\texttt{docs/PLAYSTYLE\_DATA\_COLLECTION\_PROTOCOL.md}도 추가했다. 실제 데이터셋에서는
\texttt{annotations[]}의 다중 annotator 라벨을 majority vote로 집계하고 pairwise agreement를
함께 보고한다. 또한 저장소에 있는 공개 샘플을 바탕으로
\texttt{datasets/validation/playstyle\_annotation\_packet.json}과
\texttt{datasets/validation/real\_playstyle\_sessions\_unlabeled.json}을 생성했다.
annotation packet은 라벨러 편향을 줄이기 위해 모델 예측을 포함하지 않는다.

\section{결과}
\subsection{테스트 실행 결과}
\begin{table}[h]
\centering
\begin{tabular}{lrrl}
\toprule
항목 & 통과 & 전체 & 비고 \\
\midrule
백엔드 pytest & 36 & 36 & coverage 60.29\% \\
프론트엔드 Vitest & 31 & 31 & 5개 test file 통과 \\
최종 검증 스크립트 & 15 & 15 & PASS \\
\bottomrule
\end{tabular}
\caption{1단계 및 최종 검증 결과}
\end{table}

\subsection{기존 종합 평가 재현}
\begin{table}[h]
\centering
\begin{tabular}{ll}
\toprule
지표 & 결과 \\
\midrule
파싱 정확도 & 100.00\% \\
성능 요약 & 99,109 events/sec \\
일관성 & 통과 \\
엣지 케이스 처리 & 4/4 \\
실제 데이터 파일 & 2개 \\
\bottomrule
\end{tabular}
\caption{\texttt{comprehensive\_test.py} 재실행 요약}
\end{table}

이전 분석에서 발견된 계획 시간 반례는 수정됐다. 이벤트가
\texttt{inventory\_change(1000ms)}, \texttt{inventory\_change(2000ms)},
\texttt{block\_place(5000ms)}일 때 현재 구현은 첫 준비 행동부터 첫 건축까지의 4000ms를
반환한다. 이 정의는 평가 스크립트와 회귀 테스트에 고정했으며, 이벤트 순서가 뒤섞여도 timestamp
기준의 첫 건축 이벤트를 사용한다.

\subsection{무작위 반례 탐색}
\begin{table}[h]
\centering
\begin{tabular}{lrrr}
\toprule
변형 & trial & 성공 & 실패 \\
\midrule
drop\_timestamp & 40 & 40 & 0 \\
drop\_position & 40 & 40 & 0 \\
unknown\_type & 40 & 40 & 0 \\
negative\_timestamp & 40 & 40 & 0 \\
shuffle\_order & 40 & 40 & 0 \\
\bottomrule
\end{tabular}
\caption{고정 seed 기반 무작위 강건성 탐색 결과}
\end{table}

총 200/200 trial이 성공했다. 다만 이 결과는 파서의 예외 안전성과 기본 불변식 유지에 대한
증거일 뿐, 실제 사용자 행동 분포나 성격 추론 타당도를 입증하지는 않는다.

\subsection{성능 벤치마크}
\begin{table}[h]
\centering
\begin{tabular}{rrrr}
\toprule
이벤트 수 & 반복 & median ms & p95 ms \\
\midrule
10 & 30 & 0.052 & 0.066 \\
50 & 30 & 0.169 & 0.185 \\
100 & 30 & 0.346 & 0.368 \\
500 & 30 & 2.610 & 2.690 \\
1000 & 30 & 7.683 & 7.834 \\
5000 & 30 & 140.361 & 148.367 \\
\bottomrule
\end{tabular}
\caption{이벤트 수 단일 변수 성능 실험}
\end{table}

시각화 자료는 \texttt{docs/figures/empirical\_performance.svg}와
\texttt{docs/figures/empirical\_robustness.svg}에 생성했다. 첫 번째 그림은 이벤트 수 증가에
따른 median/p95 latency를 보여주며, 두 번째 그림은 무작위 단일 변형별 성공률을 보여준다.

\subsection{플레이 스타일 카테고리 결과}
\begin{table}[h]
\centering
\begin{tabular}{lr}
\toprule
지표 & 결과 \\
\midrule
Synthetic labeled sessions & 6 \\
Correct primary category predictions & 6 \\
Fixture accuracy & 100.00\% \\
Failures & 0 \\
\bottomrule
\end{tabular}
\caption{\texttt{validate\_playstyle\_categories.py} 결과}
\end{table}

이 결과는 카테고리 scoring mechanics의 재현성 검증이다. 실제 사용자 플레이 스타일 일반화
성능을 주장하려면 사람이 라벨링한 실제 세션 데이터가 필요하다. 검증 스크립트는 실제 데이터가
제공될 경우 annotator majority label, confusion table, failure examples, pairwise agreement를
동일한 JSON 형식으로 산출한다.

추가로, 공개 Minecraft 샘플을 라벨링 전 데이터셋으로 변환해 검증 스크립트에 통과시켰다.
이 경우 라벨이 없으므로 accuracy는 \texttt{null}이며, 이는 의도된 동작이다. 라벨러가
\texttt{annotations[]}를 채운 뒤 같은 명령을 실행하면 실제 라벨 기준 정확도와 합의도가 산출된다.

\section{공개 데이터 점검}
저장소에 포함된 공개 데이터는 \texttt{opendota\_real\_match\_8650963582.json}
(converted event 12개)와 \texttt{full\_pipeline\_result.json}(raw event 8개)로 확인됐다.
두 파일 모두 파이프라인 예시로는 유용하지만, 표본 수가 작아 일반화 성능을 주장하기에는
부족하다. 또한 README에 명시된 TESS와 workout video 대용량 데이터는 현재 저장소에 없으므로
오디오/모션 기반 분석은 이번 실증분석 범위에서 제외했다.

\section{분석과 리스크}
\begin{itemize}
  \item 실행 가능성은 개선 후 확인됐다. 백엔드와 프론트엔드 테스트가 모두 통과한다.
  \item 파서의 기본 예외 안전성은 무작위 반례 탐색에서 양호했다.
  \item 성능은 1000 이벤트 p95 7.834ms, 5000 이벤트 p95 148.367ms로, 일반적인 수백~천 단위 세션에는 충분해 보인다.
  \item 계획 시간 정의는 첫 준비 행동부터 첫 건축 행동까지로 고정됐고 회귀 테스트가 추가됐다.
  \item synthetic playstyle fixture 6개는 모두 올바르게 분류됐다. 그러나 이는 실제 사용자 일반화 성능이 아니라 scoring mechanics 검증이다.
  \item 실제 데이터 표본이 매우 작고, 생체신호 대용량 데이터가 없어 사용자 수준 외적 타당도는 아직 입증되지 않았다.
  \item MuMax3 바이너리는 현재 환경에 없어 pre-computed pattern으로 대체됐다. 물리 시뮬레이션 성능 주장은 별도 환경에서 검증해야 한다.
\end{itemize}

\section{결론}
현재 프로젝트는 테스트와 기본 파이프라인 실행성 측면에서 재현 가능한 상태로 정리되었다.
그러나 ``실제 사용자 플레이 스타일을 일반화해서 정확히 분류한다''는 강한 주장은 아직 데이터와
평가 설계가 부족하다. 이번 실증분석에서 확실히 말할 수 있는 결론은 다음이다:
(1) 수정 후 구현은 백엔드/프론트엔드 테스트를 모두 통과한다,
(2) 공개 샘플과 무작위 단일 변형 입력에서 파이프라인은 안정적으로 실행된다,
(3) 계획 시간 지표의 평가/구현 정의 충돌은 해소됐다,
(4) synthetic fixture 기준 플레이 스타일 카테고리 scoring은 재현 가능하지만, 실제 라벨 데이터 검증은 아직 필요하다.

\section{재현 명령}
\begin{verbatim}
python3 -m pytest backend/tests/ -v
npm run test
PYTHONPATH=backend python3 backend/comprehensive_test.py
PYTHONPATH=backend python3 backend/final_verification_test.py
PYTHONPATH=backend python3 backend/empirical_analysis.py
PYTHONPATH=backend python3 backend/validate_playstyle_categories.py
PYTHONPATH=backend python3 backend/prepare_playstyle_annotation_packet.py
\end{verbatim}

\end{document}
